% Exploratory analyses X1a, X1b and X2, integrated 2026-09-28/29.
% Sources: detection-boundary-v4-exploratory/ declarations, REPORT.md and
% results/x1a.json, x1b.json and x2.json. EX01--EX13 are integrated ledger
% rows; review C corrections (2026-09-29) map to RCxx rows. The numbers of
% this file are transcribed from those results files, which lie outside the
% values catalog, and are kept unchanged. Terminology and the description of
% the spline basis were revised during the journal revision (agent S,
% 2026-10-13). No confirmatory result is reported here.

\subsection{Mean-field twin and variance explained by influence (X1a, X1b, X2)}\label{sec:x1-x2}\phantomsection\label{supp:x1-x2}

\paragraph{Status.}
Both analyses were declared on 28 September 2026, before their implementation and before any confirmatory result, in answer to the internal pre-results review; they leave the frozen v4 protocol unchanged and cannot be promoted to confirmatory evidence. % EX01 DS10
They used seeds disjoint from every confirmatory namespace, read no gate, tier or $a^*(r)$ output, and ran at Git revision \texttt{2f1777bb} on an exploration host with hash-verified frozen sources. % EX02
A declared second stage, X1b, estimated the mean-field twin's own boundary against B, $a^*_{MF}(r)$, with the frozen fitting, validation, band, censoring and inversion rules of the main curve, after a successor pipeline reproduced the frozen pipeline byte for byte in four cells with the original generator. % EX03
All nine X1b units executed all nine strengths and passed the zero-strength check; this operational status does not imply a gap-free band. % EX13 RC01
Four cells were gaps: $a=0.60$ at $r=0.50$, and $a=0.90$ at $r=0.55$, 0.75 and 0.80. % EX13 RC02
The first and sustained boundaries in these four units were censored at the gaps. % EX13 RC03
The gap-free grids at $r=0.60$, 0.65, 0.70 and 0.85 did not attain AUC 0.70 in their point estimates. % EX13 RC04
At $r=0.90$, the first and sustained point boundaries were $a=0.30$, with AUC 0.706 and boundary intervals $[0.30,>0.90]$; the upper end was not reached on the grid. % EX13 RC05
No lower band exceeded 0.70 at this $r$, so the point attainment does not support a strict one-sided crossing claim. % EX13 RC06
The comparison of $a^*_{MF}(r)$ with the confirmatory $a^*(r)$ was made only after the results became available; it is shown side by side and labeled exploratory (Table~\ref{tab:v4x-x1b-boundaries}). % EX13 EX12
The analyses below neither estimate nor imply the confirmatory $a^*(r)$. % EX12

\paragraph{X1a: mean-field twin.}
X1a asks whether the fixed discriminator separates the original generator from a Gaussian mean-field twin with matched conditional drift mean and variance. % EX04 RC07
Its twin is the frozen influence arm with the drift replaced by $m_t(z_{it})+\zeta_{it}\sqrt{v_t(z_{it})/8}$, where $m_t(z)$ and $v_t(z)$ are the exact mean and variance of $f(z_{jt}-z)$ over the current population and $\zeta_{it}$ is independent standard normal noise, renewed at every step, before the same standardization and update. % EX04
The contrast changes the persistence of the contact term, the dependence between persons and the conditional distribution of the drift. % EX04 RC08
In each of the 72 cells with $a\in G\setminus\{0\}$ and $r\in R$, both arms passed through the observed channel of the main curve, the frozen discriminant was trained on 2,000 panels per arm, and its AUC was computed on 2,000 reference networks with two events per arm, with a 95\% percentile bootstrap over the reference networks (1,000 resamples). % EX05
The reading is per cell: an interval excluding 0.5 supports separability of these two generators under the trained discriminant; an interval containing 0.5 does not establish equivalence, and there is no multiplicity correction or whole-curve claim. % EX05 RC09

Every interval contained 0.5 for $a\leq0.003$ (AUC 0.492--0.511), six of nine excluded it at $a=0.01$ (AUC 0.513--0.523 in those cells), and all excluded it from $a=0.03$ on, with AUCs of 0.534--0.565 at $a=0.03$, 0.623--0.716 at $a=0.10$, 0.840--0.974 at $a=0.30$ and at least 0.986 for $a\geq0.60$; none lay entirely below 0.5 (Table~\ref{tab:supp-x1a}). % EX06
Intervals excluded 0.5 in 51 of the 72 cells, all with $a\geq0.01$, and contained it in the other 21: the 18 cells with $a\leq0.003$ and the cells $r=0.50$, 0.60 and 0.75 at $a=0.01$. % EX06 RC10
Over 4,000 reference panels per arm, mean lag correlations of arm and twin differed by less than 0.0007 for $a\leq0.01$. % EX07
For $a\geq0.03$ the arm-minus-twin difference was positive at all five lags in 44 of 45 cells and rose with the lag in 43, for instance from 0.0002 at lag one to 0.0011 at lag five at $(a,r)=(0.03,0.70)$ and from 0.014 to 0.046 at $(0.90,0.70)$. % EX07
These AUCs compare the arm with its twin, not with a persistence model, and bear on no confirmatory boundary. % EX12

\paragraph{X2: variance explained by influence.}
X2 simulated latent paths without channel or mask ($N=10{,}626$, 11 waves) through a reimplementation that exposes the standardized drift $d_t$ (the standardized influence term $d_{it}$ of Section~\ref{sec:odd}, Submodel~3); all 360 paths, 10 replicates in each of 36 conditions, matched the frozen influence-arm path bit for bit. % EX08
Per replicate, pooled least squares over persons and update steps gave X2a, the partial $R^2$ of $d_t$ added to $[1,z_t]$ in predicting $z_{t+1}$ (promised in the v1 protocol), and X2b, its incremental $R^2$ beyond a natural cubic spline of $z_t$. % EX08
The spline is the one fixed by the declaration: a natural cubic spline in $z_t$ with five knots placed at the 0.05, 0.275, 0.5, 0.725 and 0.95 quantiles of the replicate's pooled $z_t$, linear beyond the outer knots, written in the truncated-power basis $1$, $z$ and $d_k(z)-d_4(z)$ for $k=1,2,3$, with $d_k(z)=[(z-t_k)_+^3-(z-t_5)_+^3]/(t_5-t_k)$; it has five columns, an intercept and four degrees of freedom, and is fitted by ordinary least squares together with $d_t$.
The conditions crossed $a\in G$, $b\in\{0.25,0\}$ and one or 24 steps per wave, and no reading rule was declared beyond summaries across replicates. % EX08
X2a increased with $a$ for every $b$ and step count, from 0.0005 at $a=0.001$ to 0.6607 at $a=0.90$ with $b=0.25$ and one step and from 0.0007 to 0.8647 with $b=0$ (Table~\ref{tab:supp-x2}); 24 steps changed it little except with $b=0.25$ and $a\geq0.30$, where it was lower (0.5443 at $a=0.90$). % EX09
X2b reached at most 0.0589 with $b=0.25$ and 0.3057 with $b=0$. % EX09
The two measures have different denominators: X2a divides the gain in fit by the residual variance after $[1,z_t]$, X2b by the total variance of $z_{t+1}$, which persistence dominates ($R^2$ of $[1,z_t]$ 0.62--0.94). % EX10
The spline added at most 0.025 to that $R^2$, and the drift's gain beyond the spline, divided by the residual variance as in X2a, was 0.95--1.23 times X2a (ratio of condition means, $a>0$). The specified spline of the person's own state therefore absorbed little of the additional variation attributable to the drift; the result concerns this five-column spline and does not extend to every nonlinear function of the own state. % EX10

\paragraph{Limitations.}
The X1a intervals are conditional on the trained discriminant and omit training variability. % EX11
The twin matches the drift's conditional mean and variance for a fresh draw of contacts but replaces a mean of eight kernel values by Gaussian noise. % EX11 RC11
X1a therefore does not identify a network-specific contribution, and the lag pattern does not decompose the distributional and dependence differences. % EX11 RC12
Cells with the same $r$ share seeds, as in the frozen plan, and are not independent replicates. % EX11
The per-step, latent $R^2$ values of X2 translate neither into AUCs nor to the observed scale. % EX11
Both analyses ran on one host, and cross-host numerical equivalence was not tested. % EX11

\begin{table}[htbp]
\centering
\footnotesize
\setlength{\tabcolsep}{3pt}
\caption{Mean-field twin (X1a): influence arm against its mean-field twin under the fixed discriminant, by influence strength $a$, over the nine values of $r\in R$. AUC: minimum--maximum over $r$, rounded to three decimals (1.000 means at least 0.9995). Excl.: cells whose 95\% bootstrap interval excludes 0.5. $\Delta L_k$: arm-minus-twin difference in mean lag-$k$ correlation over 4,000 reference panels per arm, $k=1,\ldots,5$; $\max|\Delta L_k|$ is taken over $k$ and $r$, and the last two columns count cells with $\Delta L_k>0$ at every $k$ and with $\Delta L_k$ rising in $k$. Values transcribed from the X1a results, which are archived as a sanitized copy (host and path fields removed) in the folder \texttt{archive/exploratory} next to the journal sources (Section~\ref{sec:rule-record}). Exploratory; these are not detection boundaries.}
\label{tab:supp-x1a}
\begin{tabular}{@{}lccccc@{}}
\toprule
$a$ & AUC & Excl. & $\max|\Delta L_k|$ & $\Delta L_k>0$ & Rising \\
\midrule
0.001 & 0.492--0.509 & 0/9 & 0.00024 & 2/9 & 1/9 \\
0.003 & 0.496--0.511 & 0/9 & 0.00032 & 3/9 & 3/9 \\
0.01 & 0.505--0.523 & 6/9 & 0.00069 & 6/9 & 3/9 \\
0.03 & 0.534--0.565 & 9/9 & 0.00177 & 8/9 & 7/9 \\
0.10 & 0.623--0.716 & 9/9 & 0.00563 & 9/9 & 9/9 \\
0.30 & 0.840--0.974 & 9/9 & 0.01764 & 9/9 & 9/9 \\
0.60 & 0.986--1.000 & 9/9 & 0.03916 & 9/9 & 9/9 \\
0.90 & 1.000--1.000 & 9/9 & 0.06514 & 9/9 & 9/9 \\
\botrule
\end{tabular}
\end{table}

\begin{table*}[htbp]
\centering
\footnotesize
\caption{Variance explained by influence (X2): realized contribution of the standardized drift to the next latent state, by influence strength $a$, repulsion $b$ and steps per wave $s$. Entries are means over 10 replicates ($N=10{,}626$, 11 waves), rounded to four decimals; standard deviations across replicates are at most 0.0024 for X2a and 0.0012 for X2b, and both measures are below $2\times10^{-5}$ at $a=0$. X2a is the partial $R^2$ of the drift beyond $[1,z_t]$ (residual-variance denominator) and X2b its incremental $R^2$ beyond the five-column natural cubic spline of $z_t$ described in the text (total-variance denominator), so the two blocks are not on one scale. Values transcribed from the X2 results, which are archived as a sanitized copy (host and path fields removed) in the folder \texttt{archive/exploratory} next to the journal sources (Section~\ref{sec:rule-record}). Exploratory.}
\label{tab:supp-x2}
\begin{tabular}{@{}lcccccccc@{}}
\toprule
 & \multicolumn{4}{c}{X2a, partial $R^2$} & \multicolumn{4}{c}{X2b, incremental $R^2$} \\
 & \multicolumn{2}{c}{$b=0.25$} & \multicolumn{2}{c}{$b=0$} & \multicolumn{2}{c}{$b=0.25$} & \multicolumn{2}{c}{$b=0$} \\
$a$ & $s=1$ & $s=24$ & $s=1$ & $s=24$ & $s=1$ & $s=24$ & $s=1$ & $s=24$ \\
\midrule
0 & 0.0000 & 0.0000 & 0.0000 & 0.0000 & 0.0000 & 0.0000 & 0.0000 & 0.0000 \\
0.001 & 0.0005 & 0.0006 & 0.0007 & 0.0007 & 0.0001 & 0.0001 & 0.0002 & 0.0002 \\
0.003 & 0.0017 & 0.0017 & 0.0021 & 0.0021 & 0.0004 & 0.0004 & 0.0006 & 0.0006 \\
0.01 & 0.0055 & 0.0055 & 0.0071 & 0.0072 & 0.0014 & 0.0014 & 0.0020 & 0.0020 \\
0.03 & 0.0164 & 0.0163 & 0.0216 & 0.0219 & 0.0040 & 0.0039 & 0.0061 & 0.0062 \\
0.10 & 0.0540 & 0.0530 & 0.0747 & 0.0753 & 0.0118 & 0.0116 & 0.0222 & 0.0224 \\
0.30 & 0.1670 & 0.1563 & 0.2417 & 0.2421 & 0.0303 & 0.0279 & 0.0771 & 0.0775 \\
0.60 & 0.3684 & 0.3261 & 0.5288 & 0.5287 & 0.0506 & 0.0427 & 0.1796 & 0.1800 \\
0.90 & 0.6607 & 0.5443 & 0.8647 & 0.8649 & 0.0589 & 0.0415 & 0.3055 & 0.3057 \\
\botrule
\end{tabular}
\end{table*}

\clearpage
